\section{The spectral pair created by slow translation}
\label{sec:r11-spectral}

The velocity--momentum identities determine the inertia of the charge Hessian. We now identify the pair of spectral values that emerges
from the critical scaling mode, including its leading coefficient. The
result assumes real-profile nondegeneracy; this hypothesis is not supplied
by phase coercivity alone.

Let $a$ be positive, elliptic, real-even and homogeneous of degree $m>1$.
Consider
\[
 iu_t=a(D)u-|u|^{2m/N}u,
 \qquad
 a(D)U+\omega U+iv\cdot\nabla U-|U|^{2m/N}U=0.
\]
All Hilbert spaces in this section are real unless explicitly complexified.
Let $Q=U_{1,0}$ be a real sharp ground wave and set
\[
 \Lambda Q=\frac N2Q+x\cdot\nabla Q,\qquad
 L_-=a(D)+1-|Q|^{2m/N},\qquad
 L_+=a(D)+1-\left(1+\frac{2m}{N}\right)|Q|^{2m/N}.
\]
We impose the following hypotheses.
\begin{enumerate}
\item The functions $\partial_jQ$ and $\Lambda Q$ have the regularity and
spatial decay needed to differentiate translations and critical scaling
in $L^2$ and in the profile equation. In particular,
$L_+\Lambda Q=-mQ$ holds in the operator sense.
\item $L_-\geq0$, $\ker L_-=\mathbb RQ$, and $L_-^{-1}$ is bounded
on $Q^\perp$. Moreover
$\ker L_+=\operatorname{span}\{\partial_jQ:1\leq j\leq N\}$,
and $L_+$ has a bounded inverse on that kernel's orthogonal complement.
\item There is a $C^1$ local profile branch $U_{\omega,v}$, closed under
critical scaling up to phase and translation. Its reduced action
$d(\omega,v)=E(U_{\omega,v})+\omega M(U_{\omega,v})/2
-v\cdot P(U_{\omega,v})/2$ is $C^2$ and satisfies
$d_\omega=M/2$ and $d_v=-P/2$. Profile differentiation is valid, and in
a translation gauge
\[
 D_{v_j}U_{1,0}=-iL_-^{-1}\partial_jQ.
\]
\item On the real space $\mathcal H=L^2(\mathbb R^N;\mathbb C)$ the
real action Hessians $\mathcal L_v$ are selfadjoint on a common dense
domain. Their kernels are exactly
\[
 K_v=\operatorname{span}_{\mathbb R}
       \{iU_{1,v},\partial_1U_{1,v},\ldots,\partial_NU_{1,v}\}.
\]
These kernel bases vary continuously in $\mathcal H$. The inverse of
$\mathcal L_v$ on $K_v^\perp$, extended by zero on $K_v$, exists and is
norm-continuous as a bounded operator on $\mathcal H$ for small $v$.
\end{enumerate}
The last hypothesis explicitly includes a uniform complementary spectral
gap. Common-domain operator continuity, or suitable norm-resolvent
continuity with an isolated kernel of fixed dimension, provides standard
ways to verify it. It is not being inferred for every dispersion symbol.

Define the real matrix and vector
\[
 T_{ij}=\langle\partial_iQ,L_-^{-1}\partial_jQ\rangle,
 \qquad c_j=\langle\Lambda Q,L_-^{-1}\partial_jQ\rangle,
\]
and the scalar
\begin{equation}\label{eq:r11-spectral-D}
 \mathcal D_Q=
 \langle\Lambda Q,L_-^{-1}\Lambda Q\rangle-c^TT^{-1}c.
\end{equation}
The inverse of $L_-$ always acts on $Q^\perp$.

\begin{theorem}[Small-velocity spectral asymptotics]
\label{thm:r11-spectral-pair}
Under the four hypotheses above, $T$ is positive definite and
$\mathcal D_Q>0$. If $m\ne2$, then for every sufficiently small nonzero
$v$, the nonzero spectrum of
\[
 \mathcal A_v=J\mathcal L_v,\qquad J=-i,
\]
in a fixed neighborhood of zero consists of exactly one simple pair
$\lambda_\pm(v)$, satisfying
\begin{equation}\label{eq:r11-lambda-square}
 \lambda_\pm(v)^2
  =\frac{(m-2)v^TTv}{\mathcal D_Q}+o(|v|^2).
\end{equation}
For $m>2$ this pair is real and
\[
 \lambda_\pm(v)=
 \pm\sqrt{\frac{(m-2)v^TTv}{\mathcal D_Q}}+o(|v|).
\]
In particular, the wave is spectrally unstable. For $1<m<2$ the pair is
purely imaginary, with
\[
 \lambda_\pm(v)=
 \pm i\sqrt{\frac{(2-m)v^TTv}{\mathcal D_Q}}+o(|v|).
\]
The latter assertion concerns the pair near zero and does not, by itself,
exclude unstable spectrum elsewhere.
\end{theorem}

\begin{proof}
Phase coercivity gives $T>0$: a nonzero directional derivative of a
nonzero $L^2$ function cannot vanish. Put
\[
 \Lambda_\perp Q=
 \Lambda Q-\sum_j(T^{-1}c)_j\partial_jQ.
\]
Then
\[
 \mathcal D_Q=
 \langle\Lambda_\perp Q,L_-^{-1}\Lambda_\perp Q\rangle\geq0.
\]
Equality would imply
$(x-\zeta)\cdot\nabla Q+NQ/2=0$ for some $\zeta$.
Thus $Q$ would be fixed by the mass-preserving dilation group about
$\zeta$. A nonzero function homogeneous of degree $-N/2$ cannot belong
to $L^2$: its radial mass integral is proportional to $\int_0^\infty dr/r$.
Hence $\mathcal D_Q>0$. The same Schur-complement formula shows that
$\mathcal D_Q$ is unchanged if a linear combination of translations is
added to $\Lambda Q$, so it is independent of the spatial center.

Write $U_v=U_{1,v}$ and
$g_0=iU_v$, $g_j=\partial_jU_v$. Their span $K_v$ is isotropic:
\begin{equation}\label{eq:r11-isotropic-kernel}
 \langle g_i,Jg_j\rangle=0.
\end{equation}
For the phase--translation pairings this follows from integration of a
mass derivative. For two translations it follows from commutation of
translation charges, or directly from the Fourier representation of
$\langle\partial_iU,-i\partial_jU\rangle_{\mathbb R}$.

Let $\Pi_v$ project orthogonally onto $K_v$, let $\Pi_c=I-\Pi_v$, and
write $\mathcal L_c=\mathcal L_v|_{K_v^\perp}$ and
$J_c=\Pi_cJ\Pi_c$. The operator
\[
 R_v(\lambda)=(\mathcal L_c+\lambda J_c)^{-1}
\]
is extended by zero on $K_v$. A Neumann series and hypothesis 4 show
that this inverse is analytic for $|\lambda|<\rho$, with a common
$\rho>0$ for small $v$. The eigenvalue problem is considered on the
complexification of $\mathcal H$; real pairings below extend bilinearly.

The equation $J\mathcal L_v\psi=\lambda\psi$ is equivalent to
$(\mathcal L_v+\lambda J)\psi=0$. Decompose $\psi=k+w$ with
$k\in K_v$ and $w\in K_v^\perp$. For $\lambda\ne0$ the complementary
equation gives
\[
 w=-\lambda R_v(\lambda)Jk.
\]
By \eqref{eq:r11-isotropic-kernel}, the kernel equation is therefore
exactly $D(\lambda,v)z=0$, where
\begin{equation}\label{eq:r11-reduced-matrix}
 D_{ij}(\lambda,v)
  =\langle g_i,J R_v(\lambda)Jg_j\rangle,
 \qquad k=\sum_{j=0}^Nz_jg_j.
\end{equation}
The common factor suppressed in this equation is $-\lambda^2$.
For $0<|\lambda|<\rho$, analytic block elimination therefore preserves
the eigenvalues and their algebraic multiplicities. Selfadjointness of
$\mathcal L_c$ and skew-adjointness of $J$ imply
\begin{equation}\label{eq:r11-reduced-transpose}
 D(-\lambda,v)=D(\lambda,v)^T.
\end{equation}

Profile differentiation yields
\[
 \mathcal L_v\partial_\omega U=-U,
 \qquad
 \mathcal L_v\partial_{v_j}U=-i\partial_jU=Jg_j.
\]
Components of the parameter derivatives in $K_v$ do not contribute,
by \eqref{eq:r11-isotropic-kernel}. Consequently
\[
 D(0,v)=C d''(1,v)C,
 \qquad C=\operatorname{diag}(1,-I_N).
\]
The lower-right block is $D_{vv}(0,v)=-B(1,v)/2$, where
$B=D_vP$. The first derivative of the branch gives
\[
 P(U_v)=2Tv+o(|v|),\qquad B(1,v)=2T+o(1).
\]
Thus the lower-right block stays invertible for small $\lambda,v$.
Take its scalar Schur complement $f(\lambda,v)$. Identity
\eqref{eq:r11-reduced-transpose} makes $f$ even in $\lambda$.
It is consequently analytic in $z=\lambda^2$ in a common disk, with
coefficients continuous in $v$.

The exact scaling identity for the reduced action gives
\begin{align}
 f(0,v)
 &=\frac{P^TB^{-1}P-(m-1)v\cdot P}{2m^2}\notag\\
 &=\frac{2-m}{m^2}v^TTv+o(|v|^2).
 \label{eq:r11-f-zero}
\end{align}
For completeness, the identity follows by differentiating the scaling
laws for the charges. With $H=d''$ and
$D_*=(m\omega,(m-1)v)^T$, they give
$HD_*=(0,-P/2)^T$ and $H_{vv}=-B/2$. Eliminating the velocity block
shows that its scalar Schur complement is
$[P^TB^{-1}P-(m-1)v\cdot P]/(2m^2\omega^2)$.
This uses no differentiated second-order remainder for the profile.

It remains to compute the coefficient of $\lambda^2$ at $v=0$.
Let $R_0$ be the $L^2$-orthogonal projection of $\Lambda Q/m$ off the
translation kernel, and put $Z_j=L_-^{-1}\partial_jQ$.
At $v=0$ the two real blocks give
\[
 \mathcal L_0=\operatorname{diag}(L_+,L_-),
 \qquad J(h,b)=(b,-h),
\]
\[
 \mathcal L_c^{-1}Jg_0=(-R_0,0),
 \qquad \mathcal L_c^{-1}Jg_j=(0,-Z_j).
\]
Expanding the exact resolvent in \eqref{eq:r11-reduced-matrix}, and
writing $b_j=\langle R_0,Z_j\rangle$, gives
\begin{align*}
 D_{00}(\lambda,0)
  &=\lambda^2\langle R_0,L_-^{-1}R_0\rangle+O(\lambda^4),\\
 D_{0v}(\lambda,0)&=\lambda b^T+O(\lambda^2),\\
 D_{v0}(\lambda,0)&=-\lambda b+O(\lambda^2),\\
 D_{vv}(\lambda,0)&=-T+O(\lambda).
\end{align*}
The scalar Schur complement is even. Its coefficient is therefore
\[
 \partial_z f(0,0)=
 \langle R_0,L_-^{-1}R_0\rangle-b^TT^{-1}b
 =\frac{\mathcal D_Q}{m^2}>0.
\]
The final equality uses invariance of the Schur expression under
addition of translations to the scaling derivative. Analyticity with
a common radius and continuity in $v$ now give
\[
 f(\lambda,v)=\frac{2-m}{m^2}v^TTv+o(|v|^2)
 +\lambda^2\left(\frac{\mathcal D_Q}{m^2}+o(1)\right)
 +O(\lambda^4),
\]
where the last remainder is uniform for small $v$.

The analytic scalar equation, regarded as an equation in $z=\lambda^2$,
has one simple real root near zero, by
$\partial_zf(0,0)>0$ and uniform analytic inversion. That root obeys
\eqref{eq:r11-lambda-square}. Its sign is the sign of $m-2$, since
$T>0$, and it is nonzero for nonzero sufficiently small $v$.
Its two square roots are simple because
$\partial_\lambda f=2\lambda\partial_zf\ne0$ there.
The invertible complementary and velocity blocks exclude further
nonzero spectrum in a sufficiently small fixed neighborhood.
\end{proof}

\begin{remark}[Center, frequency, and linear time scale]
If $Q(-x)=Q(x)$, then $c=0$ by parity. Without this symmetry, omitting
$c^TT^{-1}c$ gives an incorrect coefficient in general.
Critical scaling also gives the exact spectral covariance
\[
 \lambda(\omega,v)=
 \omega\lambda\bigl(1,v\omega^{-(m-1)/m}\bigr).
\]
Thus for $m>2$ and fixed $\omega>0$,
\[
 \lambda_+(\omega,v)=\omega^{1/m}
 \sqrt{\frac{(m-2)v^TTv}{\mathcal D_Q}}+o(|v|).
\]
The reciprocal growth rate is of order $|v|^{-1}$. This appendix alone concerns linear exponential growth; Section~\ref{sec:r12-nonlinear} proves a nonlinear exit time in the specified quartic class. When $m=2$, the displayed leading coefficient vanishes; a general
homogeneous degree-two symbol need not be a quadratic form, so an exact
Galilean degeneracy is not asserted.
\end{remark}

\begin{corollary}[Critical Hartree interactions]
Theorem~\ref{thm:r11-spectral-pair} holds, with the same coefficient,
for the critical Hartree flow
\[
 iu_t=a(D)u-(K*|u|^2)u,
 \qquad K(x)=|x|^{-m},\quad 1<m<N,
\]
whenever its real ground wave and local branch satisfy the same four
hypotheses with
\[
 L_-=a(D)+1-(K*Q^2),\qquad
 L_+h=L_-h-2Q\bigl(K*(Qh)\bigr).
\]
It also holds for a positive-type anisotropic interaction of the same
homogeneity under those operator and regularity assumptions.
\end{corollary}
\begin{proof}
The proof above uses the nonlinearity through phase and translation
invariance, conjugation symmetry, critical scaling, and the stated
real-block operators. These properties hold for the Hartree energy.
In particular $L_+\Lambda Q=-mQ$, the parameter equations and the
charge-scaling identities are unchanged. No step in the resolvent
reduction uses locality of $L_+$. Thus the proof applies verbatim.
\end{proof}

\begin{remark}[Relation to established instability theory]
The general mechanism of a spectral pair leaving zero at charge
Hessian degeneracy belongs to established Hamiltonian bifurcation
theory. Here the additional calculation is its coefficient for
mass-critical translation branches, including the coupling to the
translation modes. The theorem proves spectral instability for
$m>2$ under its explicit assumptions. Nonlinear orbital instability
requires a sufficiently differentiable local flow and the associated
linear evolution bounds or unstable-manifold construction in a
specified phase space; it is not inferred here for every admissible
symbol and dimension. Likewise, real-profile nondegeneracy is not
proved by this section.
\end{remark}
